\documentclass{article}
\usepackage{amsmath,amssymb}
\usepackage[vocab={actuarial}]{leanify}
\newtheorem{lemma}{Lemma}
\begin{document}

\section{Notation}

We write $S_{B,c}$ for the survival function of the Gompertz law with parameters $B$ and $c$.
We write $\mu_{B,c}$ for the force of mortality of the Gompertz law with parameters $B$ and $c$.
We write $e(\mu)$ for the life expectancy with force of mortality $\mu$.
We write $A(\mu,\delta)$ for the net single premium of a whole life insurance with force of mortality $\mu$ and force of interest $\delta$.
We write $a(\mu,\delta)$ for the present value of a continuous whole life annuity with force of mortality $\mu$ and force of interest $\delta$.
We write $E(S)$ for the complete expectation of life for the survival function $S$.
For every real number $t$, let $X(t)\coloneqq\exp(-|t|)$.

\section{Constant force}

\begin{lemma}\label{constant}
$e(\frac12)=2$, $A(1,1)=\frac12$ and $a(1,1)=\frac12$.
\end{lemma}

\begin{lemma}\label{identity}
Let $\mu,\delta>0$ be real numbers. Then $A(\mu,\delta)+\delta a(\mu,\delta)=1$.
\end{lemma}

\section{Survival models}

\begin{lemma}\label{gompertz}
$S_{1,2}(0)=1$ and $\mu_{1,2}(1)=2$, and $E(X)=1$.
\end{lemma}

\end{document}
